\chapter{公理・定義公理・基礎定理の一覧}\label{app:registry}

以下は, 証明生成器(\texttt{prover/theory.js}, \texttt{prover/found.js})に登録され, 証明木の葉として用いられる論理式の一覧である. 略記は第\ref{ch:logic}章の通り.

\section{ZFC の公理}

外延性公理 \textsf{Ext} と, 分出公理の図式 \textsf{Sep} の各インスタンス(第\ref{ch:zfc}章)を用いる. \textsf{Sep} のインスタンスであることは, 生成器が「$\forall z_1 \cdots \forall a\,\exists y\,\forall x\,(x \in y \leftrightarrow x \in a \land \varphi)$ で $y$ は $\varphi$ に自由に現れない」という形を検査して確かめる.
\begin{registrylist}
\item[Ext] $\forall x\,\forall y\,(\forall z\,(z \in x \leftrightarrow z \in y) \to x = y)$
\end{registrylist}


\section{定義公理}

それぞれ第\ref{ch:definitions}章の意味での定義による拡大の定義公理である. 関数記号の定義公理については, 対応する $\exists !$ の証明を本文の該当箇所で与えた.
\begin{registrylist}
\item[Def$\subseteq$] $\forall x\,\forall y\,(x \subseteq y \leftrightarrow \forall u \in x\,(u \in y))$
\item[Def$<$] $\forall x\,\forall y\,(x < y \leftrightarrow x \subseteq y \land x \neq y)$
\item[Def$\emptyset$] $\forall u\,(u \notin \emptyset)$
\item[Def$\cup$] $\forall x\,\forall y\,\forall u\,(u \in x \cup y \leftrightarrow u \in x \lor u \in y)$
\item[Def$\bigcup$] $\forall X\,\forall u\,(u \in \textstyle\bigcup X \leftrightarrow \exists Y \in X\,(u \in Y))$
\item[Def$S$] $\forall x\,\forall u\,(u \in S(x) \leftrightarrow u \in x \lor u = x)$
\item[DefInd] $\forall x\,(\mathrm{Ind}(x) \leftrightarrow \emptyset \in x \land \forall y \in x\,(S(y) \in x))$
\item[Def$\omega$] $\forall n\,(n \in \omega \leftrightarrow \forall x\,(\mathrm{Ind}(x) \to n \in x))$
\item[Def$-$] $\forall x\,\forall y\,(x - y = x + -y)$
\item[Def$\lvert\cdot\rvert$] $\forall x\,(0 \leq x \land \lvert x \rvert = x \lor \neg 0 \leq x \land \lvert x \rvert = -x)$
\item[Def$2$] $2 = 1 + 1$
\item[DefUB] $\forall a\,\forall X\,(\mathrm{UB}(a,X) \leftrightarrow a \in \mathbb{R} \land \forall x \in X\,(x \leq a))$
\item[DefLB] $\forall a\,\forall X\,(\mathrm{LB}(a,X) \leftrightarrow a \in \mathbb{R} \land \forall x \in X\,(a \leq x))$
\item[DefSup] $\forall s\,\forall X\,(\mathrm{IsSup}(s,X) \leftrightarrow \mathrm{UB}(s,X) \land \forall w\,(\mathrm{UB}(w,X) \to s \leq w))$
\item[DefInf] $\forall s\,\forall X\,(\mathrm{IsInf}(s,X) \leftrightarrow \mathrm{LB}(s,X) \land \forall w\,(\mathrm{LB}(w,X) \to w \leq s))$
\item[DefBddAb] $\forall X\,(\mathrm{BddAb}(X) \leftrightarrow \exists a\,(\mathrm{UB}(a,X)))$
\item[DefBddBe] $\forall X\,(\mathrm{BddBe}(X) \leftrightarrow \exists a\,(\mathrm{LB}(a,X)))$
\item[Def$X_a$] $\forall a\,\forall y\,(y \in X_{a} \leftrightarrow y \in \mathbb{R} \land \exists n \in \mathbb{N}\,(y = \underline{n} \cdot a))$
\item[Def$\ominus$] $\forall X\,\forall y\,(y \in \ominus X \leftrightarrow y \in \mathbb{R} \land -y \in X)$
\item[DefConv] $\forall a\,\forall \alpha\,(\mathrm{Conv}(a,\alpha) \leftrightarrow \mathrm{Seq}(a) \land \alpha \in \mathbb{R} \land \forall \varepsilon \in \mathbb{R}\,(0 < \varepsilon \to \exists N \in \mathbb{N}\,\forall n \in \mathbb{N}\,(N \leq n \to \lvert a_{n} - \alpha \rvert < \varepsilon)))$
\item[DefCgt] $\forall a\,(\mathrm{Cgt}(a) \leftrightarrow \exists \alpha\,(\mathrm{Conv}(a,\alpha)))$
\item[Def$\mathrm{DvgP}$] $\forall a\,(\mathrm{DvgP}(a) \leftrightarrow \mathrm{Seq}(a) \land \forall K \in \mathbb{R}\,(0 < K \to \exists N \in \mathbb{N}\,\forall n \in \mathbb{N}\,(N \leq n \to K < a_{n})))$
\item[DefSeqBdd] $\forall a\,(\mathrm{SeqBdd}(a) \leftrightarrow \mathrm{Seq}(a) \land \exists M \in \mathbb{R}\,\forall n \in \mathbb{N}\,(\lvert a_{n} \rvert \leq M))$
\item[DefMonoInc] $\forall a\,(\mathrm{MonoInc}(a) \leftrightarrow \mathrm{Seq}(a) \land \forall n \in \mathbb{N}\,(a_{n} \leq a_{S(n)}))$
\item[DefMonoDec] $\forall a\,(\mathrm{MonoDec}(a) \leftrightarrow \mathrm{Seq}(a) \land \forall n \in \mathbb{N}\,(a_{S(n)} \leq a_{n}))$
\item[DefCauchy] $\forall a\,(\mathrm{Cauchy}(a) \leftrightarrow \mathrm{Seq}(a) \land \forall \varepsilon \in \mathbb{R}\,(0 < \varepsilon \to \exists N \in \mathbb{N}\,\forall m\,\forall n\,(m \in \mathbb{N} \to n \in \mathbb{N} \to N \leq m \to N \leq n \to \lvert a_{m} - a_{n} \rvert < \varepsilon)))$
\item[DefPNbhd] $\forall f\,\forall c\,(\mathrm{PNbhd}(f,c) \leftrightarrow \exists r\,(r \in \mathbb{R} \land 0 < r \land \forall x \in \mathbb{R}\,(0 < \lvert x - c \rvert \to \lvert x - c \rvert < r \to x \in \mathrm{dom}(f))))$
\item[DefFLim] $\forall f\,\forall c\,\forall \alpha\,(\mathrm{FLim}(f,c,\alpha) \leftrightarrow \mathrm{RFn}(f) \land c \in \mathbb{R} \land \alpha \in \mathbb{R} \land \mathrm{PNbhd}(f,c) \land \forall \varepsilon \in \mathbb{R}\,(0 < \varepsilon \to \exists \delta\,(\delta \in \mathbb{R} \land 0 < \delta \land \forall x \in \mathrm{dom}(f)\,(0 < \lvert x - c \rvert \to \lvert x - c \rvert < \delta \to \lvert f(x) - \alpha \rvert < \varepsilon))))$
\item[Def$\mathbb{N}$] $\forall n\,(n \in \mathbb{N} \leftrightarrow n \in \omega \land n \neq \emptyset)$
\item[Def$1_\omega$] $1_{\omega} = S(\emptyset)$
\item[DefSubIdx] $\forall s\,(\mathrm{SubIdx}(s) \leftrightarrow \mathrm{NSeq}(s) \land \forall k \in \mathbb{N}\,(s_{k} < s_{S(k)}))$
\item[DefSubSeq] $\forall b\,\forall a\,(\mathrm{SubSeq}(b,a) \leftrightarrow \mathrm{Seq}(b) \land \exists s\,(\mathrm{SubIdx}(s) \land \forall k \in \mathbb{N}\,(b_{k} = a_{s_{k}})))$
\item[DefRNbhd] $\forall f\,\forall c\,(\mathrm{RNbhd}(f,c) \leftrightarrow \exists r\,(r \in \mathbb{R} \land 0 < r \land \forall x \in \mathbb{R}\,(c < x \to x < c + r \to x \in \mathrm{dom}(f))))$
\item[DefLNbhd] $\forall f\,\forall c\,(\mathrm{LNbhd}(f,c) \leftrightarrow \exists r\,(r \in \mathbb{R} \land 0 < r \land \forall x \in \mathbb{R}\,(c - r < x \to x < c \to x \in \mathrm{dom}(f))))$
\item[DefFLimR] $\forall f\,\forall c\,\forall \alpha\,(\mathrm{FLimR}(f,c,\alpha) \leftrightarrow \mathrm{RFn}(f) \land c \in \mathbb{R} \land \alpha \in \mathbb{R} \land \mathrm{RNbhd}(f,c) \land \forall \varepsilon \in \mathbb{R}\,(0 < \varepsilon \to \exists \delta\,(\delta \in \mathbb{R} \land 0 < \delta \land \forall x \in \mathrm{dom}(f)\,(c < x \to x < c + \delta \to \lvert f(x) - \alpha \rvert < \varepsilon))))$
\item[DefFLimL] $\forall f\,\forall c\,\forall \alpha\,(\mathrm{FLimL}(f,c,\alpha) \leftrightarrow \mathrm{RFn}(f) \land c \in \mathbb{R} \land \alpha \in \mathbb{R} \land \mathrm{LNbhd}(f,c) \land \forall \varepsilon \in \mathbb{R}\,(0 < \varepsilon \to \exists \delta\,(\delta \in \mathbb{R} \land 0 < \delta \land \forall x \in \mathrm{dom}(f)\,(c - \delta < x \to x < c \to \lvert f(x) - \alpha \rvert < \varepsilon))))$
\item[DefFLimInf] $\forall f\,\forall \alpha\,(\mathrm{FLimInf}(f,\alpha) \leftrightarrow \mathrm{RFn}(f) \land \alpha \in \mathbb{R} \land \exists r \in \mathbb{R}\,\forall x \in \mathbb{R}\,(r < x \to x \in \mathrm{dom}(f)) \land \forall \varepsilon \in \mathbb{R}\,(0 < \varepsilon \to \exists L \in \mathbb{R}\,\forall x \in \mathrm{dom}(f)\,(L < x \to \lvert f(x) - \alpha \rvert < \varepsilon)))$
\item[DefContIn] $\forall f\,\forall A\,\forall c\,(\mathrm{ContIn}(f,A,c) \leftrightarrow \mathrm{RFn}(f) \land A \subseteq \mathrm{dom}(f) \land c \in A \land \forall \varepsilon \in \mathbb{R}\,(0 < \varepsilon \to \exists \delta\,(\delta \in \mathbb{R} \land 0 < \delta \land \forall x \in A\,(\lvert x - c \rvert < \delta \to \lvert f(x) - f(c) \rvert < \varepsilon))))$
\item[DefCont] $\forall f\,\forall c\,(\mathrm{Cont}(f,c) \leftrightarrow \mathrm{ContIn}(f,\mathrm{dom}(f),c))$
\item[DefContOn] $\forall f\,\forall A\,(\mathrm{ContOn}(f,A) \leftrightarrow \mathrm{RFn}(f) \land A \subseteq \mathrm{dom}(f) \land \forall c \in A\,(\mathrm{ContIn}(f,A,c)))$
\item[DefRCont] $\forall f\,\forall c\,(\mathrm{RCont}(f,c) \leftrightarrow \mathrm{RFn}(f) \land c \in \mathrm{dom}(f) \land \forall \varepsilon \in \mathbb{R}\,(0 < \varepsilon \to \exists \delta\,(\delta \in \mathbb{R} \land 0 < \delta \land \forall x \in \mathrm{dom}(f)\,(c \leq x \to x < c + \delta \to \lvert f(x) - f(c) \rvert < \varepsilon))))$
\item[DefLCont] $\forall f\,\forall c\,(\mathrm{LCont}(f,c) \leftrightarrow \mathrm{RFn}(f) \land c \in \mathrm{dom}(f) \land \forall \varepsilon \in \mathbb{R}\,(0 < \varepsilon \to \exists \delta\,(\delta \in \mathbb{R} \land 0 < \delta \land \forall x \in \mathrm{dom}(f)\,(c - \delta < x \to x \leq c \to \lvert f(x) - f(c) \rvert < \varepsilon))))$
\item[Def{$[a,b]$}] $\forall a\,\forall b\,\forall x\,(x \in [a,b] \leftrightarrow x \in \mathbb{R} \land a \leq x \land x \leq b)$
\item[DefUC] $\forall f\,\forall A\,(\mathrm{UC}(f,A) \leftrightarrow \mathrm{RFn}(f) \land A \subseteq \mathrm{dom}(f) \land \forall \varepsilon \in \mathbb{R}\,(0 < \varepsilon \to \exists \delta\,(\delta \in \mathbb{R} \land 0 < \delta \land \forall x\,\forall y\,(x \in A \to y \in A \to \lvert x - y \rvert < \delta \to \lvert f(x) - f(y) \rvert < \varepsilon))))$
\item[DefStrInc] $\forall f\,\forall A\,(\mathrm{StrInc}(f,A) \leftrightarrow \forall x\,\forall y\,(x \in A \to y \in A \to x < y \to f(x) < f(y)))$
\item[DefItv] $\forall I\,(\mathrm{Itv}(I) \leftrightarrow I \subseteq \mathbb{R} \land \forall x\,\forall y\,\forall z\,(x \in I \to z \in I \to y \in \mathbb{R} \to x \leq y \to y \leq z \to y \in I))$
\item[DefInj] $\forall f\,(\mathrm{Inj}(f) \leftrightarrow \forall x\,\forall y\,(x \in \mathrm{dom}(f) \to y \in \mathrm{dom}(f) \to f(x) = f(y) \to x = y))$
\end{registrylist}


\section{基礎定理}\label{app:found}

次は, 本文(主に第\ref{part:formal}部)で\textbf{通常の数学の証明}を与えた $\ZFC$ の定理であり, 完全展開した証明木はまだ作成していない. 本ノートの完全展開した証明木はすべて, これらを葉として用いる. 従って, 本ノートの証明木が示しているのは「$\ZFC$ とこれらの基礎定理から, 公理・MP・Gen のみで各定理が導かれる」ことである. これらの完全展開は今後の課題である.
\begin{itemize}
  \item (R1)--(R16): $\R$ が完備順序体であること(第\ref{ch:reals}章 定理\ref{thm:complete-ordered-field}. ただし (R11)(R12) は [O1][O2] として完全展開済み).
  \item [N3]--[N6], [N9]--[N11], [N17], [N18]: 自然数 $\omega$, $\N$ の順序と埋め込み $\iota$ の性質(第\ref{ch:naturals}章 定理\ref{thm:omega-order}, 第\ref{ch:ofield}章).
  \item [S1]--[S23]: 関数値と, 定義 $(x \in A \mapsto t)$ で作った数列・関数(項ごとの演算, 部分列, 関数の和・積・商・合成など)の性質(第\ref{ch:zfc}章). [S15] は偶数番目・奇数番目の添字列の性質(第\ref{ch:naturals}章の $\omega$ の演算から).
  \item [S24], [S25]: 累乗の列 $r^\bullet$ と部分和の列 $\Sigma a$ の存在と再帰式(第\ref{ch:naturals}章 定理\ref{thm:recursion-omega}).
  \item [S26]: 単射の逆関数の性質(第\ref{ch:zfc}章の関数の基本性質から).
  \item [S27]: 指数関数の存在(第\ref{ch:text4}章 初等関数の節. 冪級数による構成).
  \item [S28]: 部分列の抽出($\omega$ の整列性と再帰定理による. 第\ref{ch:text3}章).
  \item [CC1]--[CC6]: 可算選択の図式「$\forall k \in \N\,\exists x \in \R\,\psi(k, x) \to \exists q\,(\Seq(q) \land \forall k \in \N\,\psi(k, q_k))$」のインスタンス(選択公理による. 第\ref{ch:text4}章 定理1.7 の前の説明).
\end{itemize}
\begin{registrylist}
\item[(R1)] $0 \in \mathbb{R}$
\item[(R1)] $1 \in \mathbb{R}$
\item[(R1)] $\forall x\,\forall y\,(x \in \mathbb{R} \to y \in \mathbb{R} \to x + y \in \mathbb{R})$
\item[(R1)] $\forall x\,\forall y\,(x \in \mathbb{R} \to y \in \mathbb{R} \to x \cdot y \in \mathbb{R})$
\item[(R1)] $\forall x \in \mathbb{R}\,(-x \in \mathbb{R})$
\item[(R1)] $\forall x \in \mathbb{R}\,(x \neq 0 \to x^{-1} \in \mathbb{R})$
\item[(R2)] $\forall x\,\forall y\,\forall z\,(x \in \mathbb{R} \to y \in \mathbb{R} \to z \in \mathbb{R} \to x + y + z = x + (y + z))$
\item[(R3)] $\forall x\,\forall y\,(x \in \mathbb{R} \to y \in \mathbb{R} \to x + y = y + x)$
\item[(R4)] $\forall x \in \mathbb{R}\,(x + 0 = x)$
\item[(R5)] $\forall x \in \mathbb{R}\,(x + -x = 0)$
\item[(R6)] $\forall x\,\forall y\,\forall z\,(x \in \mathbb{R} \to y \in \mathbb{R} \to z \in \mathbb{R} \to x \cdot y \cdot z = x \cdot (y \cdot z))$
\item[(R7)] $\forall x\,\forall y\,(x \in \mathbb{R} \to y \in \mathbb{R} \to x \cdot y = y \cdot x)$
\item[(R8)] $\forall x \in \mathbb{R}\,(x \cdot 1 = x)$
\item[(R8)] $0 \neq 1$
\item[(R9)] $\forall x \in \mathbb{R}\,(x \neq 0 \to x \cdot x^{-1} = 1)$
\item[(R10)] $\forall x\,\forall y\,\forall z\,(x \in \mathbb{R} \to y \in \mathbb{R} \to z \in \mathbb{R} \to x \cdot (y + z) = x \cdot y + x \cdot z)$
\item[(R13)] $\forall x\,\forall y\,(x \in \mathbb{R} \to y \in \mathbb{R} \to x < y \lor x = y \lor y < x)$
\item[(R14)] $\forall x\,\forall y\,\forall z\,(x \in \mathbb{R} \to y \in \mathbb{R} \to z \in \mathbb{R} \to x < y \to x + z < y + z)$
\item[(R15)] $\forall x\,\forall y\,(x \in \mathbb{R} \to y \in \mathbb{R} \to 0 < x \to 0 < y \to 0 < x \cdot y)$
\item[(R16)] $\forall X\,(X \subseteq \mathbb{R} \to X \neq \emptyset \to \exists b\,(\mathrm{UB}(b,X)) \to \exists s\,(\mathrm{IsSup}(s,X)))$
\item[[N3]] $\forall m\,\forall n\,(m \in \mathbb{N} \to n \in \mathbb{N} \to m \cup n \in \mathbb{N})$ \hfill{\footnotesize(補題(最大値))}
\item[[N4]] $\forall n \in \mathbb{N}\,(\underline{n} \in \mathbb{R})$ \hfill{\footnotesize(補題(ι))}
\item[[N5]] $\underline{1_{\omega}} = 1$ \hfill{\footnotesize(補題(ι))}
\item[[N6]] $\forall n \in \mathbb{N}\,(\underline{S(n)} = \underline{n} + 1)$ \hfill{\footnotesize(補題(ι))}
\item[[S1]] $\forall a\,\forall n\,(\mathrm{Seq}(a) \to n \in \mathbb{N} \to a_{n} \in \mathbb{R})$ \hfill{\footnotesize(補題(関数値))}
\item[[S2]] $\forall a\,\forall b\,(\mathrm{Seq}(a) \to \mathrm{Seq}(b) \to \mathrm{Seq}(a \oplus b) \land \forall n \in \mathbb{N}\,(a \oplus b(n) = a_{n} + b_{n}))$ \hfill{\footnotesize(定義 (x∈A↦t))}
\item[[S3]] $\forall c\,\forall a\,(c \in \mathbb{R} \to \mathrm{Seq}(a) \to \mathrm{Seq}(c \odot a) \land \forall n \in \mathbb{N}\,(c \odot a(n) = c \cdot a_{n}))$ \hfill{\footnotesize(定義 (x∈A↦t))}
\item[[S4]] $\forall a\,\forall b\,(\mathrm{Seq}(a) \to \mathrm{Seq}(b) \to \mathrm{Seq}(a \otimes b) \land \forall n \in \mathbb{N}\,(a \otimes b(n) = a_{n} \cdot b_{n}))$ \hfill{\footnotesize(定義 (x∈A↦t))}
\item[[S5]] $\forall a\,(\mathrm{Seq}(a) \to \forall n \in \mathbb{N}\,(a_{n} \neq 0) \to \mathrm{Seq}(a^{\ast}) \land \forall n \in \mathbb{N}\,(a^{\ast}(n) = a_{n}^{-1}))$ \hfill{\footnotesize(定義 (x∈A↦t))}
\item[[S6]] $\forall f\,\forall x\,(\mathrm{RFn}(f) \to x \in \mathrm{dom}(f) \to f(x) \in \mathbb{R})$ \hfill{\footnotesize(補題(関数値))}
\item[[S7]] $\forall f\,\forall x\,(\mathrm{RFn}(f) \to x \in \mathrm{dom}(f) \to x \in \mathbb{R})$ \hfill{\footnotesize(補題(関数値))}
\item[[S8]] $\forall f\,\forall a\,(\mathrm{RFn}(f) \to \mathrm{Seq}(a) \to \forall n \in \mathbb{N}\,(a_{n} \in \mathrm{dom}(f)) \to \mathrm{Seq}(f \circ a) \land \forall n \in \mathbb{N}\,(f \circ a(n) = f(a_{n})))$ \hfill{\footnotesize(定義 (x∈A↦t))}
\item[[S9]] $\forall g\,\forall f\,(\mathrm{RFn}(g) \to \mathrm{RFn}(f) \to \mathrm{RFn}(g \circ f) \land \forall x \in \mathrm{dom}(g \circ f)\,(x \in \mathrm{dom}(f) \land f(x) \in \mathrm{dom}(g) \land g \circ f(x) = g(f(x))))$ \hfill{\footnotesize(定義 (x∈A↦t))}
\item[[N9]] $\forall m\,\forall n\,(m \in \omega \to n \in \omega \to m \leq S(n) \to m \leq n \lor m = S(n))$ \hfill{\footnotesize(定理(ωの順序))}
\item[[N10]] $\forall m\,\forall n\,(m \in \omega \to n \in \omega \to m < n \to S(m) \leq n)$ \hfill{\footnotesize(定理(ωの順序))}
\item[[N11]] $\forall m\,\forall n\,(m \in \omega \to n \in \omega \to m \leq n \lor n < m)$ \hfill{\footnotesize(定理(ωの順序))}
\item[[S10]] $\forall c \in \mathbb{R}\,(\mathrm{Seq}(\overline{c}) \land \forall n \in \mathbb{N}\,(\overline{c}(n) = c))$ \hfill{\footnotesize(定義 (x∈A↦t))}
\item[[S11]] $\forall a\,(\mathrm{Seq}(a) \to \mathrm{Seq}(\lvert a \rvert) \land \forall n \in \mathbb{N}\,(\lvert a \rvert(n) = \lvert a_{n} \rvert))$ \hfill{\footnotesize(定義 (x∈A↦t))}
\item[[S12]] $\forall a\,\forall s\,(\mathrm{Seq}(a) \to \mathrm{NSeq}(s) \to \mathrm{Seq}(a \circ s) \land \forall k \in \mathbb{N}\,(a \circ s(k) = a_{s_{k}}))$ \hfill{\footnotesize(関数の合成)}
\item[[S13]] $\forall s\,\forall k\,(\mathrm{NSeq}(s) \to k \in \mathbb{N} \to s_{k} \in \mathbb{N})$ \hfill{\footnotesize(補題(関数値))}
\item[[S14]] $\mathrm{NSeq}(\mathsf{s}) \land \forall k \in \mathbb{N}\,(\mathsf{s}(k) = S(k))$ \hfill{\footnotesize(定義 (x∈A↦t))}
\item[[S15]] $\mathrm{NSeq}(\mathsf{e}) \land \mathrm{NSeq}(\mathsf{o}) \land \forall k \in \mathbb{N}\,(\mathsf{e}(k) < \mathsf{e}(S(k)) \land \mathsf{o}(k) < \mathsf{o}(S(k))) \land \forall n \in \mathbb{N}\,\exists k \in \mathbb{N}\,(n = \mathsf{e}(k) \lor n = \mathsf{o}(k))$ \hfill{\footnotesize(定理(ωの演算))}
\item[[N17]] $\forall n \in \omega\,(n < S(n))$ \hfill{\footnotesize(定理(ωの順序))}
\item[[S16]] $\forall a\,(\mathrm{Seq}(a) \to \mathrm{ran}(a) \subseteq \mathbb{R} \land \forall y\,(y \in \mathrm{ran}(a) \leftrightarrow \exists j \in \mathbb{N}\,(y = a_{j})))$ \hfill{\footnotesize(値域の定義)}
\item[[S28]] $\forall a\,\forall c\,(\mathrm{Seq}(a) \to \forall k \in \mathbb{N}\,\forall N \in \mathbb{N}\,\exists n\,(n \in \mathbb{N} \land N \leq n \land \lvert a_{n} - c \rvert < \underline{k}^{-1}) \to \exists s\,(\mathrm{SubIdx}(s) \land \forall k \in \mathbb{N}\,(\lvert a_{s_{k}} - c \rvert < \underline{k}^{-1})))$ \hfill{\footnotesize(定理(部分列の抽出))}
\item[[S17]] $\forall f\,\forall g\,(\mathrm{RFn}(f) \to \mathrm{RFn}(g) \to \mathrm{RFn}((f + g)) \land \forall x\,(x \in \mathrm{dom}((f + g)) \leftrightarrow x \in \mathrm{dom}(f) \land x \in \mathrm{dom}(g)) \land \forall x \in \mathrm{dom}((f + g))\,((f + g)(x) = f(x) + g(x)))$ \hfill{\footnotesize(関数の和)}
\item[[S18]] $\forall f\,\forall g\,(\mathrm{RFn}(f) \to \mathrm{RFn}(g) \to \mathrm{RFn}((f \cdot g)) \land \forall x\,(x \in \mathrm{dom}((f \cdot g)) \leftrightarrow x \in \mathrm{dom}(f) \land x \in \mathrm{dom}(g)) \land \forall x \in \mathrm{dom}((f \cdot g))\,((f \cdot g)(x) = f(x) \cdot g(x)))$ \hfill{\footnotesize(関数の積)}
\item[[S19]] $\forall l\,\forall f\,(l \in \mathbb{R} \to \mathrm{RFn}(f) \to \mathrm{RFn}((l f)) \land \forall x\,(x \in \mathrm{dom}((l f)) \leftrightarrow x \in \mathrm{dom}(f)) \land \forall x \in \mathrm{dom}((l f))\,((l f)(x) = l \cdot f(x)))$ \hfill{\footnotesize(関数の定数倍)}
\item[[S20]] $\forall f\,\forall g\,(\mathrm{RFn}(f) \to \mathrm{RFn}(g) \to \mathrm{RFn}((f / g)) \land \forall x\,(x \in \mathrm{dom}((f / g)) \leftrightarrow x \in \mathrm{dom}(f) \land x \in \mathrm{dom}(g) \land g(x) \neq 0) \land \forall x \in \mathrm{dom}((f / g))\,((f / g)(x) = f(x) \cdot g(x)^{-1}))$ \hfill{\footnotesize(関数の商)}
\item[[S21]] $\forall f\,(\mathrm{RFn}(f) \to \mathrm{RFn}(\lvert f \rvert) \land \forall x\,(x \in \mathrm{dom}(\lvert f \rvert) \leftrightarrow x \in \mathrm{dom}(f)) \land \forall x \in \mathrm{dom}(\lvert f \rvert)\,(\lvert f \rvert(x) = \lvert f(x) \rvert))$ \hfill{\footnotesize(関数の絶対値)}
\item[[S22]] $\forall c\,\forall r\,(c \in \mathbb{R} \to r \in \mathbb{R} \to \mathrm{Seq}(\pi_{c,r}) \land \forall n \in \mathbb{N}\,(\pi_{c,r}(n) = c + r \cdot (\underline{n} + 1)^{-1}))$ \hfill{\footnotesize(定義 (x∈A↦t))}
\item[[CC1]] $\forall f\,\forall c\,\forall \alpha\,\forall e_{0}\,(\forall k \in \mathbb{N}\,\exists x \in \mathbb{R}\,(x \in \mathrm{dom}(f) \land 0 < \lvert x - c \rvert \land \lvert x - c \rvert < \underline{k}^{-1} \land \neg \lvert f(x) - \alpha \rvert < e_{0}) \to \exists q\,(\mathrm{Seq}(q) \land \forall k \in \mathbb{N}\,(q(k) \in \mathrm{dom}(f) \land 0 < \lvert q(k) - c \rvert \land \lvert q(k) - c \rvert < \underline{k}^{-1} \land \neg \lvert f(q(k)) - \alpha \rvert < e_{0})))$ \hfill{\footnotesize(図式[CC]の例)}
\item[[S23]] $\forall g\,\forall f\,\forall x\,(\mathrm{RFn}(g) \to \mathrm{RFn}(f) \to x \in \mathrm{dom}(f) \to f(x) \in \mathrm{dom}(g) \to x \in \mathrm{dom}(g \circ f))$ \hfill{\footnotesize(関数の合成)}
\item[[CC2]] $\forall A\,\forall c\,\forall f\,\forall e_{0}\,(\forall k \in \mathbb{N}\,\exists x \in \mathbb{R}\,(x \in A \land \lvert x - c \rvert < \underline{k}^{-1} \land \neg \lvert f(x) - f(c) \rvert < e_{0}) \to \exists q\,(\mathrm{Seq}(q) \land \forall k \in \mathbb{N}\,(q(k) \in A \land \lvert q(k) - c \rvert < \underline{k}^{-1} \land \neg \lvert f(q(k)) - f(c) \rvert < e_{0})))$ \hfill{\footnotesize(図式[CC]の例)}
\item[[CC3]] $\forall a\,\forall b\,\forall f\,(\forall k \in \mathbb{N}\,\exists x \in \mathbb{R}\,(x \in [a,b] \land \underline{k} < f(x)) \to \exists q\,(\mathrm{Seq}(q) \land \forall k \in \mathbb{N}\,(q(k) \in [a,b] \land \underline{k} < f(q(k)))))$ \hfill{\footnotesize(図式[CC]の例)}
\item[[CC4]] $\forall a\,\forall b\,\forall \sigma\,\forall f\,(\forall k \in \mathbb{N}\,\exists x \in \mathbb{R}\,(x \in [a,b] \land \sigma - \underline{k}^{-1} < f(x)) \to \exists q\,(\mathrm{Seq}(q) \land \forall k \in \mathbb{N}\,(q(k) \in [a,b] \land \sigma - \underline{k}^{-1} < f(q(k)))))$ \hfill{\footnotesize(図式[CC]の例)}
\item[[CC5]] $\forall a\,\forall b\,\forall f\,\forall e_{0}\,(\forall k \in \mathbb{N}\,\exists x \in \mathbb{R}\,(x \in [a,b] \land \exists y\,(y \in [a,b] \land \lvert x - y \rvert < \underline{k}^{-1} \land \neg \lvert f(x) - f(y) \rvert < e_{0})) \to \exists q\,(\mathrm{Seq}(q) \land \forall k \in \mathbb{N}\,(q(k) \in [a,b] \land \exists y\,(y \in [a,b] \land \lvert q(k) - y \rvert < \underline{k}^{-1} \land \neg \lvert f(q(k)) - f(y) \rvert < e_{0}))))$ \hfill{\footnotesize(図式[CC]の例)}
\item[[CC6]] $\forall a\,\forall b\,\forall p\,\forall f\,\forall e_{0}\,(\forall k \in \mathbb{N}\,\exists x \in \mathbb{R}\,(x \in [a,b] \land \lvert p(k) - x \rvert < \underline{k}^{-1} \land \neg \lvert f(p(k)) - f(x) \rvert < e_{0}) \to \exists q\,(\mathrm{Seq}(q) \land \forall k \in \mathbb{N}\,(q(k) \in [a,b] \land \lvert p(k) - q(k) \rvert < \underline{k}^{-1} \land \neg \lvert f(p(k)) - f(q(k)) \rvert < e_{0})))$ \hfill{\footnotesize(図式[CC]の例)}
\item[[S24]] $\forall r \in \mathbb{R}\,(\mathrm{Seq}({(r)^{\bullet}}) \land {{r}^{1_{\omega}}} = r \land \forall n \in \mathbb{N}\,({{r}^{S(n)}} = {{r}^{n}} \cdot r))$ \hfill{\footnotesize(再帰的定義)}
\item[[S25]] $\forall a\,(\mathrm{Seq}(a) \to \mathrm{Seq}(\Sigma a) \land {\textstyle\sum_{k=1}^{1_{\omega}} a_{k}} = a_{1_{\omega}} \land \forall n \in \mathbb{N}\,({\textstyle\sum_{k=1}^{S(n)} a_{k}} = {\textstyle\sum_{k=1}^{n} a_{k}} + a_{S(n)}))$ \hfill{\footnotesize(再帰的定義)}
\item[[N18]] $\forall n \in \mathbb{N}\,(n \neq 1_{\omega} \to \exists m \in \mathbb{N}\,(n = S(m)))$ \hfill{\footnotesize(定理(ωの順序))}
\item[[S26]] $\forall f\,(\mathrm{RFn}(f) \to \mathrm{Inj}(f) \to \mathrm{RFn}(f^{-1}) \land \forall y\,(y \in \mathrm{dom}(f^{-1}) \leftrightarrow \exists x \in \mathrm{dom}(f)\,(y = f(x))) \land \forall x \in \mathrm{dom}(f)\,(f^{-1}(f(x)) = x) \land \forall y \in \mathrm{dom}(f^{-1})\,(f^{-1}(y) \in \mathrm{dom}(f) \land f(f^{-1}(y)) = y))$ \hfill{\footnotesize(逆関数)}
\item[[S27]] $\mathrm{RFn}(\exp) \land \mathrm{dom}(\exp) = \mathbb{R} \land \forall x \in \mathbb{R}\,(0 < \exp(x)) \land \forall x\,\forall y\,(x \in \mathbb{R} \to y \in \mathbb{R} \to \exp(x + y) = \exp(x) \cdot \exp(y)) \land \mathrm{StrInc}(\exp,\mathbb{R}) \land \mathrm{ContOn}(\exp,\mathbb{R}) \land \forall y \in \mathbb{R}\,(0 < y \to \exists x \in \mathbb{R}\,(\exp(x) = y))$ \hfill{\footnotesize(指数関数)}
\end{registrylist}

